\section{Konfigurationen von TLS}
Dieses Kapitel beschreibt ein Testszenario, das ein gehärtetes und ein verwundbares TLS-System gegenüberstellt. 

\subsection{Sichere Konfiguration}
Client-Server-Paar~1 simuliert eine moderne, sichere TLS-Konfiguration. Die wichtigsten Parameter sind Protokollversion, Schlüsselaustausch, Verschlüsselungsmodus, Integritätsschutz, Zertifikatstyp und Schlüssellänge.

Als Protokollversion wird TLS~1.3 verwendet~\cite{rfc8446}. Der Schlüsselaustausch erfolgt über ECDHE, sodass Perfect Forward Secrecy erreicht wird. Verschlüsselung und Integritätsschutz setzt das Paar über AEAD im GCM-Modus um~\cite{gcm-wiki}. Als Zertifikat dient ein CA-signiertes Zertifikat mit ECDSA oder 4096-Bit RSA~\cite{ecdsa-ref,bsi-tr02102}.

\subsection{Veraltete Konfiguration}
Das Client-Server-Paar~2 stellt ein ungepatchtes Altsystem dar und dient als Primärziel für die späteren Fehler-Simulationen. Es beinhaltet die unsichere TLS-Konfiguration.

Es nutzt TLS~1.2 mit rückwärtskompatiblen, unsicheren Mechanismen~\cite{rfc8996}. Der Schlüsselaustausch erfolgt statisch per RSA~\cite{rsa-studyflix}, wodurch keine Forward Secrecy besteht. Verschlüsselung per CBC und Integritätsschutz per SHA-1 MAC sind anfällig für Padding-Angriffe. Das Zertifikat ist selbstsigniert und verwendet nur 1024-Bit RSA~\cite{bsi-tr02102}.

\subsection{Konfiguration der Simulationen}

Die exakten OpenSSL-Befehle und das Setup sind im Anhang (Abschnitt~\ref{app:simulation-config}) dokumentiert. Dort finden Sie die Zertifikatsgenerierung, Client-Server-Setups sowie die Szenarien für erfolgreichen und nicht-erfolgreichen Downgrade.